\documentstyle[% 


righttag% 


,11pt% 


,amssymb% 


,russian%               remove in Eng. 


,izvru%                 Set izven% in Eng. 


]{amsart} 


 


%  


% \newcommand{\ker}{\operatorname{ker}} 


\newcommand{\const}{\operatorname{const}} 


\newcommand{\tts}[1]{}  


 


\theoremstyle{plain} 


\theorembodyfont{\it} 


\newtheorem{thmnonum}{\hskip\parindent Теорема} 


\renewcommand{\thethmnonum}{} 


 


 


%\hoffset=-15mm%                  For Eng. 


\setcounter{page}{8} 


\god{2002} 


\nomer{7~(482)} %                        Remove in Eng. 


%\nomer{Vol.\,46, No.\,7}%               Set in Eng. 


%\str{\pageref{begin}--\pageref{end}}%  Set in Eng. 


\udk{517.946} 


 


\author{\it В.В.\,Дубровский, Л.В.\,Смирнова} 


% NB:   ^^ remove in Eng. 


\title{К вопросу о восстановлении потенциала  \\


в обратной задаче Робина} 


 


\date{} 


%\pagestyle{myheadings}%         Set in Eng. 


\pagestyle{plain} %              Remove in Eng. 


\begin{document} 


%\thispagestyle{plain}%          Set in Eng. 


\maketitle 


\label{begin} 


 


 


В предлагаемой статье рассматривается вопрос о возможности восстановления 


потенциала по неполным спектральным данным в обратной задаче 


Борга--Левинсона 


с краевыми условиями третьего рода или, что то же самое, с краевыми 


условиями Робина. Под обратными задачами спектрального анализа понимают 


задачи восстановления оператора по его заданным спектральным 


характеристикам. 


Такими характеристиками могут быть спектры (при различных граничных 


условиях), спектральная функция, данные рассеяния и др. 


Основная идея приложений 


обратных задач заключается в следующем: измеряются определенные величины, 


которые можно измерить, и на основании этого пытаются получить информацию 


об интересующих физических величинах. Так в обратной задаче рассеяния на 


потенциале физической величиной, представляющей интерес, является 


потенциал 


$q(x)$ уравнения Шр\"едингера, а измеряемой величиной --- амплитуда 


рассеяния. 


К 


задачам такого типа приводят, в частности, некоторые проблемы квантовой 


механики, например, определение внутриатомных сил по заданным уровням 


энергии, 


т.\,е. по спектру, который может быть найден экспериментально. 


 


Первые работы в этом направлении появились в 1929 году. 


Была доказана 


 


\begin{thmnonum}[{\series{m}\shape{n}\selectfont [1]}] 


Пусть $\lambda_0<\lambda_1<\lambda_2<\dotsb$ --- собственные значения 


задачи Штурма--Лиувилля 


\begin{gather*} 


-y''+q(x)y=\lambda y,\quad 0\le x\le\pi,\\ 


y'(0)=y'(\pi)=0, 


\end{gather*} 


где $q(x)$ --- действительная непрерывная функция. Если $\lambda_n=n^2$, 


$n=0,1,2,\dotsc$, то 


$q(x)\equiv0$. 


\end{thmnonum} 


 


Таким образом, можно сказать, что впервые обратная задача была 


поставлена в [1]. В простейшей постановке она заключается 


в том, чтобы определить оператор, зная его спектр. 


 


Борг [2] выполнил систематическое 


исследование обратной задачи для классического оператора Штурма--Лиувилля 


по спектрам. Он показал, что в общем случае оператор Штурма--Лиувилля 


однозначно не определяется одним спектром, т.\,е. результат 


В.А.\,Амбарцумяна 


[1] является скорее исключением из общего правила. 


 


Затем появились результаты исследований, выполненных Левинсоном 


[3], [4]. Наиболее значимые результаты в теории обратных задач получены 


для дифференциального оператора Штурма--Лиувилля, для которого достаточно 


полно решена проблема корректности постановки, а также методов решения 


обратных задач. 


 


Теоремы о существовании потенциала в обратной задаче спектрального 


анализа известны лишь для обыкновенных дифференциальных уравнений [5]--[7] 


и др. или для степеней оператора Лапласа с потенциалом на прямоугольнике 


[8]. Кроме этого, многие обратные задачи имеют не единственное решение. 


Поэтому одной из важных становится проблема единственности восстановления 


потенциала, в решении которой возникает вопрос о выявлении дополнительных 


условий, обеспечивающих единственность решения обратной задачи. 


 


В [9] была доказана многомерная 


теорема Борга--Левинсона для задачи 


\begin{gather*} 


(-\Delta+q)u=\lambda u\ \ \text{в}\ \ \Omega, \\ 


u|_S=0, 


\end{gather*} 


где $\Omega$ --- ограниченная область в $R^n$ ($n\ge2$) с границей $S$ 


класса $C^\infty$. Пусть 


$\lambda_1<\lambda_2<\dotsb$ --- собственные значения задачи, 


взятые с учетом кратности, $m_j$ -- 


порядок $\lambda_j$, а $u_1,u_2,\dotsc,u_{m_j}$ --- действительнозначные 


ортонормальные собственные 


функции, соответствующие $\lambda_j$. После введения определенным образом 


отношения 


эквивалентности на множествах 


$$ 


E_j=\bigg\{\bigg(\frac{\partial u_1}{\partial v}, 


\frac{\partial u_2}{\partial v},\dotsc, 


\frac{\partial u_{m_j}}{\partial v}\bigg)\bigg|_S\bigg\}, 


$$ 


и обозначения классов эквивалентности через $W_j$, была доказана 


 


\begin{thmnonum}[{\series{m}\shape{n}\selectfont [9]}] Пусть $q_1$, 


$q_2$ --- действительнозначные функции из $C^\infty(\ov\Omega)$ такие, что 


\begin{alignat*}{2} 


\lambda_j(q_1)&=\lambda_j(q_2) &\quad &\text{для всех}\ \ j\ge1,\\ 


W_j(q_1)&=W_j(q_2) &\quad 


&\text{для всех}\ \ j\ge1. 


\end{alignat*} 


Тогда $q_1=q_2$.


\end{thmnonum} 


 


В [10], [11] получены подобные результаты и для других задач. 


 


Факт о том, что в теорему об устойчивости решения обратных задач входят 


величины, не влияющие на единственность восстановления потенциала, 


стал ясным после того, как была доказана 


 


\begin{thmnonum}[{\series{m}\shape{n}\selectfont [12]}]


Пусть $q_1$, $q_2$ --- действительнозначные функции из 


$C^\infty(\ov\Omega)$ такие, что 


\begin{alignat*}{2} 


\lambda_j(q_1)&=\lambda_j(q_2) &\quad &\text{для всех}\ \ j\ge N,\\ 


W_j(q_1)&=W_j(q_2) &\quad 


&\text{для всех}\ \ j\ge N, 


\end{alignat*} 


где $N$ --- достаточно большое натуральное число. 


Тогда $q_1=q_2$. 


\end{thmnonum} 


 


Проблема единственности восстановления потенциала в обратных задачах 


спектрального анализа при отсутствии {\it бесконечного} числа спектральных 


данных впервые изучалась в [13]--[16]. 


 


Вопрос о единственности восстановления потенциала в задаче Неймана в 


случае, если отсутствует {\it бесконечное} число собственных значений и 


значений 


собственных функций на границе заданной области, был поднят в 


[14], где доказано, что существует подпоследовательность собственных 


чисел краевой задачи Неймана, не влияющая на единственность восстановления 


потенциала. 


 


Несколько позднее в [15], [16] было показано, что если $\Omega$ --- 


ограниченная 


область в $R^2$ с границей $S$ класса $C^2$, то при определенном 


асимптотическом 


разложении собственных значений и при выполнении некоторых 


условий следует единственность восстановления потенциала в обратных 


задачах для задач Дирихле и Неймана. 


 


Вернемся к вопросу о возможности восстановления потенциала по неполным 


спектральным данным в обратной задаче Борга--Левинсона с краевыми 


условиями Робина. 


 


Пусть $\Omega$ --- ограниченная область в $R^n$, $n\ge2$, с границей $S$ 


класса $C^\infty$. 





Рассмотрим краевые задачи Робина для действительных функций 


$q_j\in C^\infty(\ov\Omega)$, $j=1,2$, и действительной функции 


$\delta\in C^\infty(S)$, $\delta(x)\le0$, 


\begin{equation} 


\begin{array}{c} 


-\Delta u(x)+q_j(x)u(x)=\lambda u(x),\quad x\in\Omega,\\[3mm] 


\displaystyle\bigg[\dfrac{\partial u}{\partial v}-\delta u 


\displaystyle\bigg]\bigg|_S=0, 


\end{array}


\end{equation} 


где $\lambda$ --- комплексный параметр, $v$ --- внутренняя нормаль к $S$. 


Под решением этих 


задач будем понимать функции из $C^\infty(\ov\Omega)$, удовлетворяющие 


всюду на $\ov\Omega$ 


соответственно каждой из них. 


 


Пусть $\mu_t(q_j)$ --- собственные числа оператора $-\Delta+q_j$, 


перечисленные в 


порядке возрастания их величин, взятые с учетом кратности, $u_t(q_j)$


--- 


соответствующие им собственные 


ортонормированные $L_2(\ov\Omega)$ функции. В случае кратных 


собственных чисел есть произвол в выборе ортонормированных функций. 


Введем оператор Дирихле $D:L_2(S)\to L_2(S)$, задаваемый равенством 


$D(\lambda,q_j)f=V_j|_S$, где $f\in C^\infty(\ov\Omega)$, а функции 


$v_j\in C^\infty(\ov\Omega)$, рассматриваемые как 


элементы из $L_2(\ov\Omega)$, являются решениями задач Робина


\begin{gather*} 


-\Delta V(x)+q_j(x)V(x)=\lambda V(x),\quad x\in\Omega,\\ 


\bigg[\frac{\partial V}{\partial v}-\delta V\bigg]\bigg|_S= 


\frac{\partial f}{\partial v}-\delta f,\quad 


\lambda\ne\mu_t(q_j),\quad t=\ov{1,\infty}. 


\end{gather*} 


 


Определим функцию рассеяния 


$$ 


F(\lambda,\omega,\theta;q_j;\delta)=\int_S\bigg(D(\lambda,q_j) 


\bigg(\frac{\partial\varphi_{\lambda,\omega}}{\partial v}- 


\delta\varphi_{\lambda,\omega}\bigg)\bigg)(x) 


\bigg(\ov{\frac{\partial\varphi_{\lambda,-\theta}}{\partial v} 


(x)-\delta\varphi_{\lambda,-\theta}(x) 


}\bigg)d S_x, 


$$ 


где $S_x$ --- элемент площади поверхности, 


$\varphi_{\lambda,\omega}(x)=\exp(i\sqrt{\lambda}\omega x)$, 


$\lambda\in C/(-\infty;0)$, $\omega,\theta\in S^{n-1}$, $\omega 


x=\sum\limits^n_{r=1}\omega_k x_k$. 


 


Можно показать, что 


\begin{multline*} 


F(\lambda,\omega,\theta;q_j)=-\frac{\lambda}{2}(\theta-\omega)^2 


\int_\Omega\exp(-i\sqrt{\lambda}(\theta-\omega)x)d x-\int_\Omega 


\exp(-i\sqrt{\lambda}(\theta-\omega)x)q_j(x)d x-\\ 


% 


-\int_S \delta(x)\exp(-i\sqrt{\lambda}(\theta-\omega)x)d S_x- 


\int_\Omega q_j(x)(-\Delta+q_j-\lambda)^{-1}(q_j\varphi_{\lambda,\omega} 


)(x)\ov{\varphi_{\lambda,-\theta}(x)}d x. 


\end{multline*} 


При этом если положить 


$$ 


c_N=\bigg(1-\frac{|\xi|^2}{4N^2}\bigg)^{1/2},\ \ 


\sqrt{l_N}=N+i,\ \ \theta_N=c_N\eta+\frac{\xi}{2N}, 


\ \ \omega_N=c_N\eta-\frac{\xi}{2N}, 


$$ 


где $(\xi,\eta)=0$, $0\ne\xi\in R^n$, имеем 


\begin{equation} 


\lim_{N\to\infty}F(l_N,\theta_N,\omega_N;q_j)= 


-\frac{|\xi|^2}{2}\int_\Omega\exp(-i x\xi)d x+ 


\int_\Omega\exp(-i x\xi)q_j(x)d x. 


\end{equation} 


 


В [9] утверждается, что формально оператор Дирихле имеет ядро 


$$ 


\ker D=\sum^\infty_{t=1}u_t(q_j)(x) 


\ov{u_t(q_j)(y)}(\mu_t(q_j)-\lambda)^{-1}, 


\ \ x,y\in S. 


$$ 


Если к тому же $\ker D(\lambda,q_1)-\ker D(\lambda,q_2)$ имеет смысл как 


оператор в $L_2(S)$, то 


учитывая, что константы интегрирования равны нулю, получаем 


\begin{multline*} 


\ker D(\lambda,q_1)-\ker D(\lambda,q_2)=\sum^\infty_{t=1} 


(\mu_t(q_2)-\mu_t(q_1))(\mu_t(q_1)-\lambda)^{-1}


(\mu_t(q_2)-\lambda)^{-1}\times \\ 


% 


\times u_t(q_1)(x)\ov{u_t(q_1)(y)}+\sum^\infty_{t=1} 


(u_t(q_1)(x)\ov{u_t(q_1)(y)}-u_t(q_2)(x) 


\ov{u_t(q_2)(y)})(\mu_t(q_2)-\lambda)^{-1}. 


\end{multline*} 


И тогда для операторной нормы имеем 


\begin{multline} 


\|\ker D(\lambda,q_1)-\ker D(\lambda,q_2)\|\le\sum^\infty_{t=1} 


|\mu_t(q_2)-\mu_t(q_1)|\,|\mu_t(q_1)-\lambda|^{-1} 


|\mu_t(q_2)-\lambda|^{-1}\|u_t(q_1)(x)\|^2_{L_2(S)}+ \\ 


% 


+\sum^\infty_{t=1} 


\|u_t(q_1)(x)\ov{u_t(q_1)(y)}-u_t(q_2)(x) 


\ov{u_t(q_2)(y)}\|_{L_2(S\times S)} 


|\mu_t(q_2)-\lambda|^{-1}\equiv S_1+S_2. 


\end{multline} 


 


Пусть собственные функции задач Робина (1) с потенциалами $q_1$ и $q_2$ 


совпадают на границе $S$, а собственные значения совпадают, начиная с 


некоторого номера $T$. Тогда 


$$ 


\|\ker D(\lambda,q_1)-\ker D(\lambda,q_2)\|\le\sum^T_{t=1} 


|\mu_t(q_2)-\mu_t(q_1)|\,|\mu_t(q_1)-\lambda|^{-1} 


|\mu_t(q_2)-\lambda|^{-1} 


\|u_t(q_1)(x)\|^2_{L_2(S)}. 


$$ 





Положив $\lambda=N^2-1+2N i$, имеем 


\begin{equation} 


\lim_{N\to\infty}\|D(\lambda,q_1)-D(\lambda,q_2)\|\bigg\| 


\frac{\partial\varphi_{\lambda,\omega}}{\partial\gamma}- 


\delta\varphi_{\lambda,\omega}\bigg\|_{L_2(S)}\bigg\| 


\frac{\partial\varphi_{\lambda,-\theta}}{\partial\gamma}- 


\delta\varphi_{\lambda,-\theta}\bigg\|_{L_2(S)}=0. 


\end{equation} 


Отсюда следует, что 


$\lim\limits_{N\to\infty}|F(l_N,\theta_N,\omega_N;q_1)- 


F(l_N,\theta_N,\omega_N;q_2)|=0$. Из этого заключаем, 


что преобразования Фурье функций $q_1$ и $q_2$ совпадают, а следовательно, 


и совпадают сами эти функции. В результате справедлива теорема, 


подобная той, что получена в [12] для задач такого типа, но 


с краевыми условиями Дирихле. 


 


\begin{thm} 


Пусть $q_1$, $q_2$ --- действительные функции из $C^\infty(\ov\Omega)$ 


такие, что 


 


$1.$ $\forall t\in N$ \ $u_t(q_1)|_S=u_t(q_2)|_S$, 


 


$2.$ $\exists T\in N$ \ $\forall t\ge T$ \ $\mu_t(q_1)=\mu_t(q_2)$. 


 


Тогда $q_1(x)=q_2(x)$ для любого $x\in\ov\Omega$. 


\end{thm} 


 


Рассмотрим другие условия, при которых выполняется равенство (4). Для 


этого положим 


\begin{equation} 


\lambda=N^2-1+2N i\equiv[\mu_{t_0}]+2N i. 


\end{equation} 


 


Далее, т.\,к. $\|u_t(q_j)\|_{L_2(S)}\le C_0|\mu_t(q_j)|^\zeta$, 


$0<\zeta<\frac{n}{4}$, будем рассматривать 


только те потенциалы $q_1$, $q_2$, где 


$\mu_t(q_j)=C_1t^{2/n}+o\big(t^{1/n+\gamma}\big)$, $0<\gamma<1/n$, $C_1>0$, и 


$\|u_t(q_j)\|_{L_2(S)}\le C_2t^{1/n}$, $C_2>0$, что справедливо для 


плоских областей. 


 


Вернeмся к равенству (3). 


Пусть 


\begin{align*} 


p_0&=\underset{t\to\infty}{\varlimsup}{}t^\varepsilon 


|\mu_t(q_2)-\mu_t(q_1)|<\infty \quad \text{при} 


\ \ \varepsilon>\frac{3n+4}{3n} \\


\intertext{и} 


q_0&=\underset{t\to\infty}{\varlimsup}{}t^\alpha 


\|u_t(q_1)(x)\ov{u_t(q_1)(y)}-u_t(q_2)(x) 


\ov{u_t(q_2)(y)}\|_{L_2(S\times S)}<\infty \quad 


\text{при}\ \ \alpha>\frac{2n+1}{2n}. 


\end{align*} 


Тогда, используя тот факт [17], что при $|n-t_k|>C t^\beta_k>0$, где 


$\frac{1}{n}<\beta\le1$, 


\begin{align*} 


|\mu_n(q_j)-\mu_{t_k}(q_j)|&\ge\const\cdot\max 


\{t_k,n\}^{2/n+\beta-1}>0,\ \ j=1,2, \\ 


% 


|\mu_n(q_1)-\mu_{t_k}(q_2)|&\ge\const\cdot\max 


\|t_k,n\}^{2/n+\beta-1}>0,


\end{align*} 


получим следующие оценки: 


\begin{multline*} 


S_1t^{2/n}_0\equiv\sum_{|t-t_0|<C t^\beta_0}+ 


\sum_{|t-t_0|\ge C t^\beta_0>0}\le\const\bigg( 


\frac{p_0t^{-\varepsilon}_0t^{4/n}_0t^\beta_0}{t^{2/n}_0}+ 


\sum_{|t-t_0|\ge C t^\beta_0>0}\frac{p_0t^{-\varepsilon}t^{2/n}_0} 


{t^{2\beta}}\bigg)\le \\ 


% 


\le\const p_0\bigg(\frac{1}{t^{\varepsilon-\beta-2/n}_0}+ 


\frac{1}{t^{2\beta+\varepsilon-3}_0}\bigg)\to0 


\end{multline*} 


при $t_0\to\infty$ и $\max\big\{\frac{3-\varepsilon}{2}; 


\frac{1}{n}\big\}<\beta<\min\big\{1;\varepsilon-\frac{2}{n}\big\}$. 


Проведя некоторые рассуждения, можно показать, что такие $\beta$ всегда 


существуют. 


 


Аналогично, 


\begin{multline*} 


S_2t^{2/n}_0\equiv\sum_{|t-t_0|<C t^{\Tilde\beta}_0}+ 


\sum_{|t-t_0|\ge C t^{\Tilde\beta}_0>0}\le\const\bigg( 


\frac{q_0t^{-\alpha}_0t^{\Tilde\beta}_0t^{2/n}_0}{t^{1/n}_0}+ 


\sum_{|t-t_0|\ge C t^{\Tilde\beta}_0>0}\frac{q_0t^{-\alpha}t^{2/n}_0} 


{t^{\Tilde\beta}}\bigg)\le \\ 


% 


\le\const q_0\bigg(\frac{1}{t^{-1/n+\alpha-\Tilde\beta}_0}+ 


\frac{1}{t^{\Tilde\beta+\alpha-2}_0}\bigg)\to0 


\end{multline*} 


при $t_0\to\infty$ и $\max\big\{2-\alpha;\frac{1}{n}\big\}<\wt\beta< 


\min\big\{1;-\frac{1}{n}+\alpha\big\}$. 


Такие $\wt\beta$ также существуют. 


 


Если же $N\to\infty$, то из (5) следует, что $t_0\to\infty$, т.\,е. 


предельное 


равенство (4) справедливо. В этом случае 


$$ 


\lim_{N\to\infty}|F(l_N,\theta_N,\omega_N;q_1)- 


F(l_N,\theta_N,\omega_N;q_2)|=0, 


$$ 


и, воспользовавшись равенством (2), можем утверждать, что преобразования 


Фурье для функций $q_1$ и $q_2$ совпадают. 


 


\begin{thm} 


Пусть в задачах Робина $(1)$ потенциалы $q_j$, $j=1,2$, в $\ov\Omega$ 


таковы, что выполняются условия 


 


$1.$ $\mu_t(q_j)=C_1t^{2/n}+o\big(t^{1/n+\gamma}\big)$, 


$0<\gamma<\frac{1}{n}$, $C_1>0;$ 


 


$2.$ $\|u_t(q_j)\|_{L_2(S)}\le C_2t^{1/n}$, $C_2>0;$ 


 


$3.$ $\exists\varepsilon>\frac{3n+4}{3n}$ \ 


$p_0=\underset{t\to\infty}{\varlimsup}{}t^\varepsilon|\mu_t(q_2)= 


\mu_t(q_1)|<\infty;$ 


 


$4.$ $\exists\alpha>\frac{2n+1}{2n}$ \ 


$q_0=\underset{t\to\infty}{\varlimsup}{}t^\alpha


\|u_t(q_1)(x)\ov{u_t(q_1)(y)}- 


u_t(q_2)(x)\ov{u_t(q_2)(y)}\|_{L_2(S\times S)}<\infty$. 


 


Тогда $q_1=q_2$. 


\end{thm} 


 


 


\begin{thebibliography}{99} 


 


\bibitem{1} 


Ambartsumian V. {\em \"Uber eine Frage der Eigenwerttheorie} // 


Zeitschrift f\"ur Physik. 


-- 1929. -- Bd.\,53. -- S.\,690--695. 


\bibitem{2} 


Borg G. {\em Eine Umkehrung der Sturm--Liouvillschen Eigenwert aufgabe} // 


Acta. Math. -- 1946. -- Bd.\,78. -- \No\,1. -- S.\,1--96. 


\bibitem{3} 


Levinson N. {\em The inverse Sturm--Liouville problem}. -- 


Math. Tidsskr. B, 1949. -- P.\,25--30. 


\bibitem{4} 


Levinson N. {\em On the uniqueness of the potential in a Schr\"odinger 


equation for a given asymptotic phase} // 


Danske Vid. Selsk. Mat. Fys. Medd. -- 1949. -- V.\,25. -- 


\No\,9. -- Р.\,25. 


\bibitem{5} 


Березанский Ю.М. {\em О теореме единственности в обратной задаче 


спектрального анализа для уравнения Шр\"едингера} // Тр. Моск. матем. об-ва. 


-- 1958. -- Т.\,7. -- \No\,3. -- C.\,3--62. 


\bibitem{6} 


Левитан Б.М. {\em Обратные задачи Штурма--Лиувилля}. -- М.: Наука, 


1984. -- 240 c. 


\bibitem{7} 


Рамм А.Г. {\em Многомерные обратные задачи рассеяния}. -- М.: Мир, 1994. 


-- 469 с. 


\bibitem{8} 


Дубровский В.В., Великих А.С. {\em Теорема о существовании решения 


обратной 


задачи спектрального анализа для степени оператора Лапласа} // 


Электромагнитные волны и электронные системы. -- 1998. -- \No\,5. 


-- С.\,6--9. 


\bibitem{9} 


Nachman F., Sylvester J., Uhlmann G. {\em An $n$-dimensional 


Borg--Levinson 


theorem} // J. Math. Phys. -- 1988. -- V.\,115. -- P.\,595--605. 


\bibitem{10} 


Ramm А.С. {\em Multidimensional inverse problems and completeness of the 


products of solutions to PDE} // J. Math. Anal. Appl. -- 1988. 


-- V.\,134. -- Р.\,211--253. 


\bibitem{11} 


Suzuki T. {\em Ultrahyperbolic approach to some multidimensional inverse 


problems}. -- Proc. Japan Acad. -- 1988. -- 64 p. 


\bibitem{} 


Isozaki H. {\em Some remarks on the multidimensional Borg--Levinson 


theorem} 


// J. Math. Kyoto Univ. -- 1991. -- V.\,31. -- \No\,3. -- Р.\,743--753. 


\bibitem{13} 


Дубровский В.В. {\em К устойчивости обратных задач спектрального анализа 


для 


уравнений математической физики} // Тр. Моск. матем. об-ва. -- 1994. -- 


Т.\,49. -- Bып.\,3. -- С.\,171--172. 


\bibitem{14} 


Дубровский В.В. {\em К многомерной обратной задаче спектрального анализа} 


// 


Тр. Моск. матем. об-ва. -- 1994. -- Т.\,49. -- Bып.\,3. -- С.\,229--230. 


\bibitem{15} 


Дубровский В.В. {\em Теорема о единственности решения обратных задач 


спектрального анализа} // Дифференц. уравнения. -- 1997. -- Т.\,33. 


-- \No\,3. -- С.\,421--422. 


\bibitem{16} 


Дубровский В.В. {\em Об одном неравенстве в обратных задачах спектрального 


анализа} // Дифференц. уравнения. -- 1997. -- Т.\,33. -- \No\,6. 


-- С.\,843--844. 


\bibitem{17} 


Дубровский В.В. {\em К абстрактной формуле Гельфанда--Левитана}


// УMH. -- 1991. -- T.\,46. 


-- \No\,3. -- С.\,187--188. 


 


\end{thebibliography} 


 


\vskip 2cm 


 


\postup{\it Магнитогорский государственный}{\qquad \it Поступили} 


\postup{\it университет}{\it первый вариант $17.09.1999$} 


\postup{}{\it окончательный вариант $15.10.2001$} 


 


\label{end} 


\end{document} 


